\documentclass[10pt]{amsart}
\usepackage[a4paper,margin=22mm]{geometry}
\usepackage[no-math]{fontspec}
\setmainfont{Latin Modern Roman}[SmallCapsFont={lmromancaps10-regular.otf}]
\usepackage{amsmath,amssymb,amsthm,mathtools,hyperref,xcolor}
\usepackage[all,cmtip]{xy}
\hypersetup{hidelinks}
\emergencystretch=3em
\newtheorem{theo}{Theorem}[section]
\newtheorem{lemma}[theo]{Lemma}
\newtheorem{prop}[theo]{Proposition}
\newtheorem{coro}[theo]{Corollary}
\theoremstyle{definition}
\newtheorem{defi}[theo]{Definition}
\newtheorem{rema}[theo]{Remark}
\newtheorem{exam}[theo]{Example}
\providecommand{\ie}{i.e.}
\providecommand{\eg}{e.g.}
\providecommand{\resp}{resp.}
\providecommand{\G}{\mathfrak G}
\usepackage{tikz}
\usepackage[centertableaux]{ytableau}
\usepackage[all,cmtip]{xy}
\usetikzlibrary{arrows,matrix,positioning}
\tikzset{>=stealth',
  head/.style = {fill = white, text=black},
  plaque/.style = {draw, rectangle, minimum size = 10mm}, 
  pil/.style={->,thick},
  junct/.style = {draw,circle,inner sep=0.5pt,outer sep=0pt, fill=black}
  }

\providecommand{\noopsort}[1]{}

\renewcommand{\G}{\mathfrak{G}}  % Grothendieck
\newcommand{\wG}{\mathfrak{J}}  % weak Grothendieck
\newcommand{\hG}{\mathfrak{H}}  % canonical (hook) Grothendieck
\newcommand{\dG}{\mathfrak{g}}  % dual Grothendieck
\newcommand{\dwG}{\mathfrak{j}}  % dual weak Grothendieck
\newcommand{\xx}{\mathbf{x}}
\newcommand{\HH}{\mathcal{H}}
\newcommand{\iso}{\cong}
\newcommand{\bbb}{\mathsf{b}}
\newcommand{\gp}{\mathsf{g}}

\newcommand{\abs}[1]{\lvert #1 \rvert}
\newcommand{\fsl}{\mathfrak{sl}}
\newcommand{\gl}{\mathfrak{gl}}
\newcommand{\sym}{S}
\newcommand{\ml}[1]{ \mathbf{\color{darkred} #1} } % Marked letters

\DeclareMathOperator{\wt}{wt} % weight
\DeclareMathOperator{\rd}{rd} % reading word
\DeclareMathOperator{\ssyt}{SST} % semistandard (Young) tableaux
\DeclareMathOperator{\svt}{SVT} % set-valued tableaux
\DeclareMathOperator{\mvt}{MVT} % multiset-valued tableaux
\DeclareMathOperator{\hvt}{HVT} % hook-valued tableaux
\DeclareMathOperator{\rpp}{RPP} % reverse plane partitions
\DeclareMathOperator{\vst}{VST} % valued-set tableaux
\DeclareMathOperator{\Gr}{Gr} % Grassmannian
\DeclareMathOperator{\RSK}{RSK} % RSK

\newcommand{\ZZ}{\mathbb{Z}}
\newcommand{\CC}{\mathbb{C}}

\newcommand{\mcB}{\mathcal{B}}
\newcommand{\mcF}{\mathcal{F}}
\newcommand{\mcP}{\mathcal{P}}


\newcommand{\bplus}{{\color{blue}+}}
\newcommand{\bminus}{{\color{red}-}}


% Dark red emphasis
\definecolor{darkred}{rgb}{0.7,0,0} % darkred color
\newcommand{\defn}[1]{{\color{darkred}\emph{#1}}} % emphasis of a definition

%%%%%%%%% a placer avant \begin{document}


%%%%%%%%%%%%%%%%%%%%%%%%%%%%%%%%%%%%%
\newcommand*{\mk}{\mkern -1mu}
\newcommand*{\Mk}{\mkern -2mu}
\newcommand*{\mK}{\mkern 1mu}
\newcommand*{\MK}{\mkern 2mu}

\hypersetup{urlcolor=purple, linkcolor=blue, citecolor=red}


\newcommand*{\romanenumi}{\renewcommand*{\theenumi}{\roman{enumi}}}
\newcommand*{\Romanenumi}{\renewcommand*{\theenumi}{\Roman{enumi}}}
\newcommand*{\alphenumi}{\renewcommand*{\theenumi}{\alph{enumi}}}
\newcommand*{\Alphenumi}{\renewcommand*{\theenumi}{\Alph{enumi}}}
%%%%%%%%%%%%%%%%%%%%%%%%%%%%%%%%%%



%%%%%%%%%%%%%%%%%%%%%%%%%%%%%%%%%%%%%%%%%%%
\title[Column-flagged uncrowding]{Column-flagged uncrowding of multiset-valued tableaux and weak Grothendieck Schur multiplicities}
\author{Graham Hawkes and Travis Scrimshaw}
\date{Algebraic Combinatorics3(3)(2020),727–755; DOI10.5802/alco.111}
\begin{document}
\maketitle
\setcounter{section}{1}
\section{Background}
\label{sec:background}

In this section, we give the necessary background on the crystal structure on set-valued tableaux and on (weak) symmetric/canonical Grothendieck functions.
We use English convention for partitions and tableaux.
Let $\xx = (x_1, x_2, x_3, \ldots)$ be a countable sequence of indeterminants.
Let $\fsl_n$ denote the special linear Lie algebra (\ie the simple Lie algebra of type $A_{n-1}$) over $\CC$ and $U_q(\fsl_n)$ the corresponding Drinfel'd--Jimbo quantum group.
Let $\lambda = (\lambda_1, \lambda_2, \dotsc, \lambda_{\ell})$ be a partition; a sequence of weakly decreasing positive integers.
Let $\ell(\lambda) = \ell$ denote the length of $\lambda$.

%%%%%%%%%%
\subsection{Semistandard tableaux, set-valued tableaux, and crystals}

A \defn{(semistandard) set-valued tableau of shape $\lambda$} is a filling $T$ of the boxes of Young diagram of $\lambda$ by finite nonempty sets of positive integers so that rows are weakly increasing and columns are strictly increasing in the following sense:
For every set $A$ to the left of a set $B$ in the same row, we have $\max A \leq \min B$, and for $C$ below $A$ in the same column, we have $\max A < \min C$.
A set-valued tableau is a \defn{semistandard (Young) tableau} if all sets have size $1$.
Let $\svt^n(\lambda)$ (\resp $\ssyt^n(\lambda)$) denote the set of set-valued (\resp semistandard) tableaux of shape $\lambda$ with entries at most $n$.

In~\cite{MPS18}, a $U_q(\fsl_n)$-crystal structure, in the sense of Kashiwara~\cite{K90,K91}, was given on $\svt^n(\lambda)$.
Recalling this crystal structure, we begin with the \defn{crystal operators} $e_i, f_i \colon \svt^n(\lambda) \to \svt^n(\lambda) \sqcup \{0\}$, where $i \in I := \{1, \dotsc, n-1\}$.

\begin{defi}
\label{defn:svt_crystal_ops}
Fix some $T \in \svt^n(\lambda)$ and $i \in I$.
Write $\bplus$ above each column of $T$ containing $i$ but not $i+1$, and write $\bminus$ above each column containing $i+1$ but not $i$. Next cancel signs in ordered pairs $\bminus \bplus$ until obtaining a sequence of the form $\bplus \cdots \bplus \bminus \cdots \bminus$ called the \defn{$i$-signature}.
\begin{description}
\item[\defn{$e_i T$}] If there is not a $\bminus$ in the resulting sequence, then $e_i T = 0$. Otherwise let $\bbb$ correspond to the box of the leftmost uncanceled $\bminus$. Then $e_i T$ is given by one of the following:
\begin{itemize}
\item if there exists a box $\bbb^{\leftarrow}$ immediately to the left of $\bbb$ that contains an $i+1$, then remove the $i+1$ from $\bbb^{\leftarrow}$ and add an $i$ to $\bbb$;
\item otherwise replace the $i+1$ in $\bbb$ with an $i$.
\end{itemize}

\item[\defn{$f_i T$}] If there is not a $\bplus$ in the resulting sequence, then $f_i T = 0$. Otherwise let $\bbb$ correspond to the box of the rightmost uncanceled $\bplus$. Then $f_i T$ is given by one of the following:
\begin{itemize}
\item if there exists a box $\bbb^{\rightarrow}$ immediately to the right of $\bbb$ that contains an $i$, then remove the $i$ from $\bbb^{\rightarrow}$ and add an $i+1$ to $\bbb$;
\item otherwise replace the $i$ in $\bbb$ with an $i+1$.
\end{itemize}
\end{description}
\end{defi}

For a set-valued tableau $T \in \svt^n(\lambda)$ the \defn{weight} is defined as
\begin{equation}
\label{eq:weight_defn}
\wt(T) := x_1^{m_1} x_2^{m_2} \cdots x_n^{m_n},
\end{equation}
where $m_i$ is the number of occurrences of $i$ in $T$.
Denote $\abs{T} := \sum_{i=1}^n m_i$.
Define the statistics
\[
\varepsilon_i(T) = \max \{k \in \ZZ_{\geq 0} \mid e_i^k T \neq 0\},
\qquad\quad
\varphi_i(T) = \max \{k \in \ZZ_{\geq 0} \mid f_i^k T \neq 0\}.
\]
This gives a $U_q(\fsl_n)$-crystal structure on $\svt^n(\lambda)$.\footnote{The standard references for crystal structures consider the weight as an additive group, but we consider it as a multiplicative group because it useful for defining polynomials in the sequel.}
In particular, for $T, T' \in \svt^n(\lambda)$, we have
\[
e_i T = T' \Longleftrightarrow T = f_i T'.
\]
We say a $T \in \svt^n(\lambda)$ is \defn{highest weight} if $e_i T = 0$ for all $i \in i$.
For more details on crystals, we refer the reader to~\cite{BS17,K91}.

When we restrict this crystal structure to semistandard Young tableaux of shape $\lambda$, we exactly recover the crystal $B(\lambda)$ of the irreducible highest weight $U_q(\fsl_n)$-representation of highest weight $\lambda$~\cite{K90,K91}. Furthermore, the crystal operators from Definition~\ref{defn:svt_crystal_ops} also give a crystal structure on words of length $\ell$, which we naturally equate with the tensor product $B(\Lambda_1)^{\otimes \ell}$ (for more details, we refer the reader to~\cite{BS17}).
The Lusztig involution is an involution on highest weight crystals $\ast \colon B(\lambda) \to B(\lambda)$ that sends the highest weight element to the lowest weight element and extended as a crystal isomorphism
\[
e_i(T^*) \mapsto (f_{n+1-i}T)^*,
\qquad
f_i(T^*) \mapsto (e_{n+1-i}T)^*,
\qquad
\wt(T^*) = w_0 \wt(T),
\]
where $w_0$ is the permutation that reverses all entries. We extend this functorially to tensor products by applying it to every factor and then reversing the factors. The Lusztig involution is also given by the Sch{\"u}tzenberger involution (or evacuation) on semistandard tableaux~\cite{Lenart07}.

Recall that for two $U_q(\fsl_n)$-crystals $\mcB, \mcB'$, a \defn{strict crystal morphism} $\psi \colon \mcB \to \mcB'$ is a map $\psi \colon \mcB \sqcup \{0\} \to \mcB \sqcup \{0\}$ such that
\[
\psi(0) = 0,
\qquad
\psi(e_i b) = e_i \psi(b),
\qquad
\psi(f_i b) = f_i \psi(b),
\qquad
\wt\bigl( \psi(b) \bigr) = \wt(b),
\]
where we consider $e_i 0 = 0$ and $f_i 0 = 0$.
We say $\psi$ is an \defn{embedding} (\resp \defn{isomorphism}) if $\psi^{-1}(0) = \{0\}$ (\resp $\psi$ is a bijection).
When there exists an isomorphism $\psi \colon \mcB \to \mcB'$, we say $\mcB$ is isomorphic to $\mcB'$ and denote this by $\mcB \iso \mcB'$.

For a partition $\lambda$, we will sometimes write it as $\sum_{i=1}^n c_i \Lambda_i$, where $c_i$ denotes the number of columns of height $i$.
This is the usual identification of partitions with the dominant weights associated to $\fsl_n$ using the fundamental weights.
\subsection{Canonical Grothendieck functions}

From~\cite{Buch02}, we can define a \defn{symmetric Grothendieck function} as
\[
\G_{\lambda}(\xx; \beta) := \sum_{T \in \svt^{\infty}(\lambda)} \beta^{\abs{T}-\abs{\lambda}} \wt(T),
\]
where $\abs{\lambda}$ denotes the size of $\lambda$ (\ie the number of boxes in $\lambda$).
The value $\abs{T} - \abs{\lambda}$ is the so-called \defn{excess} statistic.
When $\beta = 0$, we recover the \defn{Schur function}:
\[
s_{\lambda}(\xx) = \sum_{T \in \ssyt^{\infty}(\lambda)} \wt(T).
\]
Note that when we restrict to $n$ variables $x_1, x_2, \dotsc, x_n$, we recover the \defn{$\beta$-character} of $\svt^n(\lambda)$ and $\G_{\lambda}$ (and $s_{\lambda}$) is a polynomial (as opposed to a formal power series).
For more on Schur functions, we refer the reader to~\cite[Ch.~7]{ECII}.

The \defn{weak symmetric Grothendieck function} is defined by
\begin{equation}
\label{eq:weak_definition}
\wG_{\lambda}(\xx; \alpha) := \G_{\lambda}\left( \frac{x_1}{1-\alpha x_1}, \frac{x_2}{1-\alpha x_2}, \frac{x_3}{1-\alpha x_3}, \ldots; \alpha \right),
\end{equation}
which recovers the definition given in~\cite[Thm.~6.11]{PP16} when $\alpha = -1$ and $x_i \mapsto -x_i$ (which is for the conjugate shape of the definition in~\cite{LamPyl07}).
Indeed, following~\cite{PP16} we have
\[
\frac{x_i}{1 - \alpha x_i} = x_i + \alpha x_i^2 + \alpha^2 x_i^3 + \cdots = \sum_{k=0}^{\infty} \alpha^k x_i^{k+1},
\]
which is equivalent to allowing multisets to fill the tableaux.
More explicitly, define a \defn{(semistandard) multiset-valued tableau of shape $\lambda$} to be a filling $T$ of the boxes of $\lambda$ by finite nonempty multisets of positive integers such that rows are weakly increasing and columns are strictly increasing in the same sense as for set-valued tableaux.
Let $\mvt^n(\lambda)$ denote the set of all multiset-valued tableaux of shape $\lambda$ and max entry $n$.
Thus, we arrive at the combinatorial definition of~\cite{LamPyl07} for weak symmetric Grothendieck functions:
\[
\wG_{\lambda}(\xx; \alpha) = \sum_{T \in \mvt^{\infty}(\lambda)} \alpha^{\abs{T}-\abs{\lambda}} \wt(T).
\]
\section{Multiset-valued tableaux}
\label{sec:crystal_MVT}

In this section, we prove our first main result: there exists a $U_q(\fsl_n)$-crystal structure on $\mvt^n(\lambda)$ such that it is isomorphic to a direct sum of irreducible highest weight crystals $B(\mu)$.
After that, we discuss the relation with the usual crystal structure on semistandard Young tableaux and some consequences for stable dual Grothendieck polynomials.


%%%%%%%%%%
\subsection{Crystal structure}

In order to define the crystal structure, we start by defining the crystal operators 
\[
e_i, f_i \colon \mvt^n(\lambda) \to \mvt^n(\lambda) \sqcup \{0\}.
\]
We do so by following the signature rule as for set-valued tableaux in Definition~\ref{defn:svt_crystal_ops}, for which we need to define an appropriate reading word.
We write multisets as words for compactness.

\begin{defi}[Reading word]
Let $T \in \mvt^n(\lambda)$.
Let $C$ be a column of $T$. Define the \defn{column reading word} $\rd(C)$ by first reading the smallest entry of each box from bottom-to-top in $C$, then reading the remaining entries from smallest to largest in each box from top-to-bottom in $C$.
Define the \defn{reading word} 
\[
\rd(T) = \rd(C_1) \rd(C_2) \dotsm \rd(C_k),
\]
 where $C_1, C_2, \dotsc, C_k$ are the columns of $T$ from left-to-right.
\end{defi}

\begin{exam}
For the multiset-valued (column) tableau
\[
\ytableausetup{boxsize=2.2em}
C = \; \ytableaushort{{\ml{1}13},{\ml{4}445},{\ml{6}},{\ml{7}899}}\,,
\]
we have
\[
\rd(C) = \mathbf{\color{darkred}7641} 13 445 899.
\]
\end{exam}

\begin{defi}[Crystal operators]
\label{defn:mvt_crystal_ops}
Fix $T \in \mvt^n(\lambda)$ and $i \in I$.
Write $\bplus$ for each $i$ in $\rd(T)$ and $\bminus$ for each $i+1$ in $\rd(T)$ (ignore all other letters). 
Next, cancel signs in ordered pairs $\bminus \bplus$ until obtaining a sequence of the form $\bplus \cdots \bplus \bminus \cdots \bminus$ called the \defn{$i$-signature}.
\begin{description}
\item[\defn{$e_i T$}] If there is no $\bminus$ in the resulting sequence, then $e_i T = 0$. Otherwise let $\bbb$ correspond to the box of the leftmost uncanceled $\bminus$. Then $e_i T$ is given by one of the following:
\begin{itemize}
\item if there exists a box $\bbb^{\uparrow}$ immediately above $\bbb$ that contains an $i$, then remove $i + 1$ from $\bbb$ and add $i$ to $\bbb^{\uparrow}$;
\item otherwise replace the $i+1$ in $\bbb$ with an $i$.
\end{itemize}

\item[\defn{$f_i T$}] If there is no $\bplus$ in the resulting sequence, then $f_i T = 0$. Otherwise let $\bbb$ correspond to the box of the rightmost uncanceled $\bplus$. Then $f_i T$ is given by one of the following:
\begin{itemize}
\item if there exists a box $\bbb^{\downarrow}$ immediately below $\bbb$ that contains an $i+1$, then remove the $i$ from $\bbb$ and add an $i+1$ to $\bbb^{\downarrow}$;
\item otherwise replace the $i$ in $\bbb$ with an $i+1$.
\end{itemize}
\end{description}
\end{defi}
\setcounter{theo}{4}
First we show that the crystal operators are well-defined and satisfy the requisite properties.

\begin{lemma}
\label{lemma:fixed_reading}
Let $T \in \mvt^n(\lambda)$ and suppose $f_i T \neq 0$, then the $i$ changed to $i+1$ in $f_i T$ does not change its position in the reading word.
That is to say, we have $\rd(f_i T) = f_i \rd(T)$.
\end{lemma}

\begin{proof}
Let $T' = f_i T$, and suppose the changed $i$ is in box $\bbb$.
Clearly if $i$ becomes an $i+1$ in the same box, then it does not change its position.
Now consider when $i$ moves from $\bbb$ and becomes an $i+1$ in the box $\bbb^{\downarrow}$ immediately below $\bbb$. In this case, the $i$ changed is the rightmost $i$ in the $i$-signature of $T$, and hence the rightmost entry of $\bbb$ in the reading word (we also have $\max \bbb = i$ since $\min \bbb^{\downarrow} = i+1$). Hence, the changed $i+1$ in $T'$ must be read immediately after all entries of $\bbb^{\downarrow}$ in $T'$ have been read (consider this added $i+1$ to be the second such letter), which means it remains at the same position.
\end{proof}
\setcounter{theo}{6}
\begin{lemma}
\label{lemma:mvt_ops_well_defined}
Let $T, T' \in \mvt^n(\lambda)$. Then
\[
e_i T \in \mvt^n(\lambda) \sqcup \{0\},
\qquad
f_i T' \in \mvt^n(\lambda) \sqcup \{0\},
\qquad
e_i T = T' \Longleftrightarrow T = f_i T'.
\]
\end{lemma}

\begin{proof}
We first show $f_i T' \in \mvt^n(\lambda) \sqcup \{0\}$, and if $f_i T' = 0$, then the claim is trivially true.
Suppose $T = f_i T'$ is formed by changing an $i$ to an $i+1$ within the same box $\bbb$. To show $T \in \mvt^n(\lambda)$, we first note that we could not have an $i+1$ in the box $\bbb^{\downarrow}$ immediately below $\bbb$. So it remains to show that there does not exist an $i$ in the box $\bbb^{\rightarrow}$ immediately to the right of $\bbb$. By semistandardness of $T'$, for the box $\bbb^{\searrow}$ immediately to the right of $\bbb^{\downarrow}$, we must have $i + 1 \notin \bbb^{\searrow}$. Hence, the $i \in \bbb^{\rightarrow}$ must correspond to an unpaired $\bplus$ to the right of the $\bplus$ from $\bbb$, which is a contradiction. Thus, we have $T \in \mvt^n(\lambda)$ in this case.

Now instead suppose $T = f_i T'$ is formed by moving an $i$ from a box $\bbb$ to an $i+1$ in $\bbb^{\downarrow}$, the box immediately below $\bbb$. Since $i+1 \in \bbb^{\downarrow}$ already in $T'$, we have 
\[
T \in \mvt^n(\lambda).
\]

Next we show the claim $e_i T = T'$ if and only if $T = f_i T'$.
From Lemma~\ref{lemma:fixed_reading} (and the analogous statement for $e_i$), the $i$-signatures of $T'$ and $T$ must differ by $\bplus \leftrightarrow \bminus$ corresponding to the $i \leftrightarrow i+1$ with no other cancellations. Hence $e_i$ changes $i+1 \mapsto i$ in the same box if and only if $f_i$ changes $i \mapsto i+1$ in the same box. Therefore, we have $e_i T = T'$ if and only if $T = f_i T'$.

Finally, the claim $e_i T \in \mvt^n(\lambda)$ follows from the other two statements (this claim could also be proven directly similar to $f_i T' \in \mvt^n(\lambda)$).
\end{proof}

Next we show that the coefficient of $\alpha^a$ for a weak symmetric Grothendieck of a single column is isomorphic to a crystal of semistandard tableaux for hook shape.
Note that 
\[
\abs{T} - k = \abs{T} - \abs{\lambda}.
\]

\begin{prop}
\label{prop:column_isomorphism}
Recall that $1^k = \Lambda_k$ is a single column of height $k$.
Let
\[
\mvt^n_a(\Lambda_k) := \{ T \in \mvt^n(\Lambda_k) \mid \abs{T} - k = a \}.
\]
Then
\[
\mvt^n_a(\Lambda_k) \iso B(\Lambda_k + a\Lambda_1).
\]
\end{prop}

\begin{proof}
Let $\mu = \Lambda_k + a \Lambda_1$.
We prove the claim by constructing an explicit crystal isomorphism $\psi \colon \mvt^n_a(\Lambda_k) \to B(\mu)$.
We define $\psi$ by
\[
\ytableausetup{boxsize=2.3em}
\def\arraystretch{1.9}
\begin{array}{|@{\;}c@{\;}|}
\hline
m_1 \leq o_1 \leq \cdots \leq o_{a_1}
\\\hline
m_2 \leq o_{a_1+1} \leq \cdots \leq o_{a_2}
\\\hline
\vdots
\\\hline
m_k \leq o_{a_{k-1}+1} \leq \cdots \leq o_{a_k}
\\\hline
\end{array}
\; \mapsto \;
  \ytableaushort{{m_1}{o_1}{\cdots}{o_{a_1}}{o_{a_1\!+1}}{\cdots}{o_{a_2}}{\cdots}{o_{a_k}},{m_2},{\raisebox{-2pt}{$\vdots$}},{m_k}}\,,
\]
We note that $m_i \leq o_{a_i} < m_{i+1} \leq o_{a_i+1}$ since the column is strictly increasing. Hence, we have $\psi(T) \in \ssyt^n(\mu)$ for all $T \in \mvt^n(\Lambda_k)$.
Note also that the reading words of $T$ and $\psi(T)$ are equal and so $\psi$ is a crystal isomorphism by Lemma~\ref{lemma:fixed_reading}.
\end{proof}

\begin{theo}
\label{thm:mvt_crystal}
Let $\lambda$ be a partition.
For any $T \in \mvt^n(\lambda)$ such that $T$ is a highest weight element, the closure of $T$ under the crystal operators is isomorphic to $B(\mu)$, where $\mu = \wt(T)$.
Moreover, we have
\[
\mvt^n(\lambda) \iso \bigoplus_{\mu \supseteq \lambda} B(\mu)^{\oplus M_{\lambda}^{\mu}},
\]
where $M_{\lambda}^{\mu}$ is the number of highest weight elements of weight $\mu$ in $\mvt^n(\lambda)$, and
\[
\wG_{\lambda}(\xx; \alpha) = \sum_{\mu \supseteq \lambda} \alpha^{\abs{\mu}-\abs{\lambda}} M_{\lambda}^{\mu} s_{\mu}(\xx).
\]
\end{theo}

\begin{proof}
Note that we can consider any multiset-valued tableau as a tensor product of single column multiset-valued tableaux by the definition of the crystal operators and the reading word.
In particular, this gives a strict crystal embedding, which implies that the image is a union of connected components.
Hence the first claim follows from Proposition~\ref{prop:column_isomorphism}, Lemma~\ref{lemma:mvt_ops_well_defined}, and that the tensor product of highest weight crystals is a direct sum of highest weight crystals.
The other two claims follow immediately from the first.
\end{proof}
\subsection{Uncrowding the crystal structure}
\label{sec:uncrowding}

An \defn{increasing tableau} is a semistandard Young tableau that is also strictly increasing across its rows.
Let $\mcF^c_{\mu/\lambda}$ denote the set of increasing tableaux of shape $\mu/\lambda$ where the $i$-th column is \defn{strictly flagged} by $i$, that is to say the maximum entry in the $i$-th column is strictly less than $i$.
We call a tableau in $\mcF^c_{\mu/\lambda}$ a \defn{column flagged tableau}.
Let $T \xleftarrow{\RSK} T'$ denote the Robinson--Schensted--Knuth (RSK) insertion (see, \eg \cite{ECII} for more details on RSK) of the reading word of $T'$ into $T$.

Next, we construct an explicit crystal isomorphism
\[
\Upsilon \colon \mvt^n(\lambda) \to \bigsqcup_{\mu \supseteq \lambda} B(\mu) \times \mcF^c_{\mu/\lambda},
\]
where the crystal structure on the codomain is given by $f_i(b \times F) = (f_i b) \times F$ for all $b \times F \in B(\mu) \times \mcF^c_{\mu/\lambda}$ for any fixed $\mu$.
We call the map $\Upsilon$ \defn{uncrowding} as it is given similar to the uncrowding map for set-valued tableaux (see~\cite[Sec.~6]{Buch02},~\cite[Sec.~5]{BM12}, and~\cite[Thm.~3.12]{MPS18}; see also~\cite{RTY18}), but working column-by-column and measuring the growth of the diagram along columns.
More specifically, for any $T \in \mvt^n(\lambda)$ we define $\Upsilon(T)$ recursively starting with $b_{\lambda_1+1} \times F_{\lambda_1+1} = \emptyset \times \emptyset$. Suppose we are at step $i$ with the current state being $b_i \times F_i$, and let $C_j$ denote the $j$-th column of $T$. Construct
\[
b_{i-1} := \rd(C_{i-1}) \xleftarrow{\RSK} \rd(C_i) \xleftarrow{\RSK} \cdots \xleftarrow{\RSK} \rd(C_{\lambda_1}).\footnote{This is equivalent to RSK inserting (the reading word of) the image under the isomorphism $\psi$ from Proposition~\ref{prop:column_isomorphism} of the rightmost $i-1$ columns of $T$.}
\]
Construct $F_{i-1}$ by starting first with $F_i$ of shape $\mu_i$ but shifting the necessary elements to the right one step, so partially filling in the shape $\mu_{i-1} / \lambda_{\geq i-1}$, where
\begin{equation}
\label{eq:rightmost_column_shape}
\lambda_{\geq i-1} = \bigl( \max(\lambda_1-i+2,0), \dotsc, \max(\lambda_{\ell}-i+2,0) \bigr)
\end{equation}
is the shape of the rightmost $i-1$ columns of $\lambda$ and $\mu_{i-1}$ is the shape of $b_{i-1}$.
Then add entries in the unfilled boxes in column $j$ with entry $j-1$ until $F_{i-1}$ has been filled in. Thus, we constructed the $(i-1)$-th step $b_{i-1} \times F_{i-1}$. Repeating this for every column, the final result is $\Upsilon(T) = b_1 \times F_1$.
\setcounter{theo}{15}
\begin{exam}
\label{ex:uncrowding_big}
Applying uncrowding to
\[
\ytableausetup{boxsize=1.8em}
T = \ytableaushort{{112}{22}{256},{33}{444}{7},{568},{9}}\,,
\]
where $\rd(T) = 953112368 \; 42244 \; 7256$,
we obtain
\begin{align*}
\ytableausetup{boxsize=1.5em}
b_4 \times F_4 & = \emptyset & \times \quad & \emptyset,  % 7256 ->
\\ b_3 \times F_3 & = \ytableaushort{256,7} & \times \quad & \ytableaushort{{\cdot}{\ml{1}}{\ml{2}},{\cdot}}\,,  % 42244 ->
\allowdisplaybreaks \\ b_2 \times F_2 & = \ytableaushort{222456,447} & \times \quad & \ytableaushort{{\cdot}{\cdot}12{\ml{4}}{\ml{5}},{\cdot}{\cdot}{\ml{2}}}\,,  % 953112367 ->
\allowdisplaybreaks \\ b_1 \times F_1 & = \ytableaushort{112222456,33447,568,9} & \times \quad & \ytableaushort{{\cdot}{\cdot}{\cdot}1245{\ml{7}}{\ml{8}},{\cdot}{\cdot}{\cdot}2{\ml{4}},{\cdot}{\ml{1}}{\ml{2}},{\cdot}}\,,
\end{align*}
resulting in $\Upsilon(T) = b_1 \times F_1$.
\end{exam}
\begin{theo}
\label{thm:flagged_decomposition}
We have
\[
\mvt^n(\lambda) \iso \bigoplus_{\mu \supseteq \lambda} B(\mu)^{\oplus \abs{\mcF^c_{\mu/\lambda}}},
\]
where the isomorphism is given by the uncrowding map $\Upsilon$, and so $M_{\lambda}^{\mu} =  \abs{\mcF^c_{\mu/\lambda}}$. Moreover, we have
\[
\wG_{\lambda}(\xx; \alpha) = \sum_{\mu \supseteq \lambda} \alpha^{\abs{\mu}-\abs{\lambda}} \abs{\mcF^c_{\mu/\lambda}} s_{\mu}(\xx).
\]
\end{theo}
\begin{proof}
It is sufficient to show that the map $\Upsilon$ is a bijection as RSK insertion is a crystal isomorphism (see, \eg \cite{BS17,LLT02}). By the construction, the result of $\Upsilon$ satisfies the flagging and row strictness conditions. By the standard properties of RSK and a straightforward induction, the result of $\Upsilon$ also satisfies the column strictness condition. So the map $\Upsilon$ is well-defined. We can construct the inverse $\Upsilon^{-1}$ by conjugating by the Lusztig involution and recursively by applying reverse RSK inserting the boxes in the column flagged tableau that are maximal in their column and the extra column added. Indeed, if we are at $b_i \times F_i$, we first apply the Lusztig involution to $b_i$ to get $b_i^*$ and perform reverse RSK insertion on the boxes of $b_i^*$ on the outer corners from bottom to top, and then on the outer box in the row corresponding to a box $\bbb$ in column $j$ of $F$ such that $j-1$ is the entry in $\bbb$, doing this from left-to-right until no longer possible. That is to say we treat all of these extra entries as being the same value in the recording tableaux. This determines the $i$-th column after applying the Lusztig involution again to the resulting word. Since $\RSK(w^*)^* = \RSK(w)$, the result is $b_{i+1} \times F_{i+1}$. Thus this is the inverse procedure of $\Upsilon$.
\end{proof}
We can construct the recording tableau we use to perform $\RSK^{-1}$ in the proof of Theorem~\ref{thm:flagged_decomposition} recursively following the description in constructing $\Upsilon^{-1}(b_1 \times F_1)$. We can clearly reconstruct $F_i$ by using the entries that are maximal in each column. For the $i$-th step, suppose the $i$-th column has height $h$, we increase all of the current entries by $h+1$, then we set the rightmost unset entry in row $j$ to $j+1$. Finally, for every entry in $F_i$ that is maximal in its column in row $j$, we set an entry to be $1$ in row $j$. We repeat this until we obtain a semistandard Young tableau.

\begin{exam}
Consider $b_1 \times F_1$ from Example~\ref{ex:uncrowding_big}. Then we construct the corresponding recording tableau $Q$ as
\begin{gather*}
\ytableaushort{{\cdot}{\cdot}{\cdot}{\cdot}{\cdot}{\cdot}112,{\cdot}{\cdot}{\cdot}13,114,5}\,,
% \allowdisplaybreaks\\
\qquad
\ytableaushort{{\cdot}{\cdot}{\cdot}112445,{\cdot}1346,447,8}\,,
\allowdisplaybreaks\\
Q = \ytableaushort{112445778,34679,77{10},{11}}\,.
\end{gather*}
Performing the Lusztig involution on $b_1$ (in $\fsl_{10}$), we obtain
\[
T' = \ytableaushort{145667799,25688,388,4}\,,
\]
and then $\RSK^{-1}(T',Q)$ gives the word $w^* = 4583\;66886\;247899751$. When we apply the Lusztig involution to $w^*$ and obtain $w = 953112368\;42244\;7256$, which is precisely the reading word of $T$, and we can separate into columns based on the values of $Q$.
\end{exam}
\begin{thebibliography}{99}
\bibitem[1]{BM12} Bandlow, Jason; Morse, Jennifer Combinatorial expansions in K-theoretic bases, Electron. J. Combin., Volume 19 (2012) no. 4, Paper no. Paper 39, 27 pages.
\bibitem[3]{Buch02} Buch, Anders Skovsted A Littlewood–Richardson rule for the K-theory of Grassmannians, Acta Math., Volume 189 (2002) no. 1, pp. 37-78.
\bibitem[6]{BS17} Bump, Daniel; Schilling, Anne Crystal bases: Representations and Combinatorics, World Scientific Publishing Co. Pte. Ltd., Hackensack, NJ, 2017, xii+279 pages.
\bibitem[16]{K90} Kashiwara, Masaki Crystalizing the q-analogue of universal enveloping algebras, Comm. Math. Phys., Volume 133 (1990) no. 2, pp. 249-260.
\bibitem[17]{K91} Kashiwara, Masaki On crystal bases of the q-analogue of universal enveloping algebras, Duke Math. J., Volume 63 (1991) no. 2, pp. 465-516.
\bibitem[19]{LamPyl07} Lam, Thomas; Pylyavskyy, Pavlo Combinatorial Hopf algebras and K-homology of Grassmannians, Int. Math. Res. Not. IMRN, Volume 2007 (2007) no. 24, Paper no. Art. ID rnm125, 48 pages.
\bibitem[20]{LLT02} Lascoux, Alain; Leclerc, Bernard; Thibon, Jean-Yves The Plactic Monoid, Algebraic Combinatorics on Words (Lothaire, M., ed.), Cambridge University Press, Cambridge, 2002.
\bibitem[24]{Lenart07} Lenart, Cristian On the combinatorics of crystal graphs. I. Lusztig’s involution, Adv. Math., Volume 211 (2007) no. 1, pp. 204-243.
\bibitem[27]{MPS18} Cara Monical, Oliver Pechenik, and Travis Scrimshaw. Crystal structures for symmetric Grothendieck polynomials. Preprint, http://arxiv.org/abs/1807.03294,2018.
\bibitem[31]{PP16} Patrias, Rebecca; Pylyavskyy, Pavlo Combinatorics of K-theory via a K-theoretic Poirier–Reutenauer bialgebra, Discrete Math., Volume 339 (2016) no. 3, pp. 1095-1115.
\bibitem[36]{RTY18} Reiner, Victor; Tenner, Bridget Eileen; Yong, Alexander Poset edge densities, nearly reduced words, and barely set-valued tableaux, J. Combin. Theory Ser. A, Volume 158 (2018), pp. 66-125.
\bibitem[39]{ECII} Stanley, Richard P. Enumerative combinatorics. Vol. 2, Cambridge Studies in Advanced Mathematics, 62, Cambridge University Press, Cambridge, 1999, xii+581 pages (With a foreword by Gian-Carlo Rota and appendix 1 by Sergey Fomin).
\end{thebibliography}
\end{document}
